% !TeX root = ../tfg.tex
% !TeX encoding = utf8
%
% ============================================================================
%  EJEMPLO 4: CLIMATIZACIÓN DE UNA PISCINA CUBIERTA (2D, CONDICIONES MIXTAS)
% ----------------------------------------------------------------------------
%  FRAGMENTO A INSERTAR EN:  capitulos/04-granjaServidores.tex
%
%  UBICACIÓN EXACTA (en 04-granjaServidores.tex):
%    * Insertar justo DESPUÉS del cierre del "Ejemplo 3: Extensión a 2D
%      Frank-Kamenetskii" (es decir, después de la línea que contiene el
%      \end{itemize} que cierra la subsección "Resultados numéricos" de la
%      sección \ref{sec:ejemplo_fk2d}).
%    * Y ANTES de la línea "\section{Contraste de los Modelos}".
%
%  Concretamente, pegar este bloque entre:
%       \end{itemize}            <- final de "Ejemplo 3"
%       (AQUÍ VA ESTE FRAGMENTO)
%       \section{Contraste de los Modelos}
%
%  NOTA: al insertarlo, las secciones "Contraste de los Modelos",
%  "Estudio de la Tasa de Convergencia" y "Convergencia en Malla" se
%  renumeran automáticamente (LaTeX lo gestiona). Opcionalmente, actualizar la
%  Tabla~\ref{tab:contraste} añadiendo una cuarta columna (se propone el
%  bloque al final de este archivo, comentado).
%
%  FIGURAS NECESARIAS (capturas de la aplicación 3D interactiva):
%    * piscina_sub_super.png   : sub-solución, super-solución y solución.
%    * piscina_convergencia.png: decaimiento del gap (escala logarítmica).
%    * piscina_app.png         : captura de la aplicación 3D interactiva.
% ============================================================================

\section{Ejemplo 4: Climatización de una Piscina Cubierta (2D con Condiciones de Contorno Mixtas)}
\label{sec:ejemplo_piscina}

\subsection{Contexto y formulación}

Este último ejemplo traslada el método a un escenario físico concreto: la
\textbf{distribución estacionaria de temperatura en una piscina cubierta}. A
diferencia de los ejemplos anteriores, aquí la frontera $\partial\Omega$ no
admite un único tipo de condición de contorno, sino que se descompone en
tres tramos con comportamientos termodinámicos distintos (envolvente acristalada,
hueco de puerta y cerramiento opaco), lo que conduce a un problema con
\textbf{condiciones de contorno mixtas}. El interés del ejemplo es doble: por un
lado, muestra que la iteración monótona se adapta sin dificultad a problemas
que no son de Dirichlet puro; por otro, el término no lineal acopla dos
mecanismos físicos de naturaleza distinta (saturación por evaporación y fuentes
de calor corporales), lo que enriquece la estructura de $f(x,u)$.

El dominio $\Omega \subset \mathbb{R}^2$ representa la planta de la instalación
(una sala rectangular con un vaso de piscina centrado). La temperatura se
trabaja en una \textbf{variable normalizada}
\begin{equation}\label{eq:piscina_normalizacion}
  u(x) = \frac{T(x) - T_{\min}}{T_{p} - T_{\min}} \in [0,1],
  \qquad
  T_{\min} = 10\,^{\circ}\mathrm{C}, \quad T_{p} = 37\,^{\circ}\mathrm{C},
\end{equation}
de modo que $u = 0$ corresponde al valor más frío (el exterior en la puerta,
$T_{\min} = 10\,^{\circ}\mathrm{C}$) y $u = 1$ a la temperatura corporal
$T_{p} \approx 37\,^{\circ}\mathrm{C}$. Esta normalización es la que permite que
la no linealidad logística $u(1-u)$ esté bien definida sobre el intervalo
$[0,1]$.

La distribución estacionaria de temperatura se rige por la ecuación de
reacción--difusión
\begin{equation}\label{eq:piscina_modelo}
  -\Delta u(x) = \lambda\, u(x)\bigl(1 - u(x)\bigr)
  + \sum_{i=1}^{N_{\mathrm{pers}}} \kappa_i(x)\,\bigl(u_{p,i} - u(x)\bigr),
  \qquad x \in \Omega,
\end{equation}
donde $\lambda > 0$ y $N_{\mathrm{pers}}$ es el número de personas presentes en
la sala. El término no lineal $f(x,u)$ acopla dos mecanismos termodinámicos
locales:

\begin{itemize}
  \item \textbf{Saturación por evaporación:} el agua transfiere calor al
  ambiente, pero la tasa térmica decae conforme el aire se satura. Se modela
  mediante la función logística $\lambda u(1-u)$, donde $\lambda > 0$ es la
  intensidad del proceso. El factor $u$ representa el impulso evaporativo
  (creciente con la temperatura) y el factor $(1-u)$ la saturación: a medida que
  el ambiente se acerca al límite caliente ($u \to 1$), el intercambio se
  amortigua.

  \item \textbf{Fuentes de calor corporales:} cada persona actúa como una fuente
  local acoplada. Cada individuo $i$ posee una temperatura corporal normalizada
  $u_{p,i} \in [0,1]$ (típicamente $u_{p,i} \approx 1$, pues
  $T_{p} \approx 37\,^{\circ}\mathrm{C}$), y la transferencia es proporcional al
  gradiente térmico entre la persona y el ambiente, $u_{p,i} - u(x)$. La función
  indicadora $\kappa_i(x)$ concentra esta contribución en la posición del
  individuo $i$: $\kappa_i(x) = \kappa > 0$ si $x$ se encuentra en la celda
  ocupada por el individuo $i$, y $\kappa_i(x) = 0$ en otro caso. Nótese que este
  término desaparece cuando la temperatura ambiente iguala la corporal del
  individuo ($u = u_{p,i}$).
\end{itemize}

\subsection{Condiciones de contorno mixtas}

La frontera $\partial\Omega$ se divide en tres tramos según el comportamiento
físico de las envolventes del edificio (Figura~\ref{fig:piscina_app}):

\begin{equation}\label{eq:piscina_cc}
  \begin{cases}
    \dfrac{\partial u}{\partial n} + \alpha(x)\, u = \alpha(x)\, u_{\mathrm{ext}},
      & x \in \partial\Omega_{\mathrm{cristal}}, \\[8pt]
    u(x) = u_{\mathrm{puerta}},
      & x \in \partial\Omega_{\mathrm{puerta}}, \\[8pt]
    \dfrac{\partial u}{\partial n} = 0,
      & x \in \partial\Omega_{\mathrm{aislada}},
  \end{cases}
\end{equation}
donde
\begin{itemize}
  \item $\partial\Omega_{\mathrm{cristal}}$ es la \textbf{cristalera} (pared
  trasera): condición de Robin (tercer tipo) que modela la convección térmica
  proporcional a la diferencia con el exterior. Aquí
  $u_{\mathrm{ext}} = (T_{\mathrm{ext}} - T_{\min})/(T_{p} - T_{\min})$
  con $T_{\mathrm{ext}} = 18\,^{\circ}\mathrm{C}$, de modo que
  $u_{\mathrm{ext}} \approx 0.30$;
  \item $\partial\Omega_{\mathrm{puerta}}$ es el \textbf{hueco de la puerta}
  automática: condición de Dirichlet pura que impone una severa pérdida de calor
  cuando la puerta se abre, forzando la temperatura local al valor exterior
  $u_{\mathrm{puerta}} = 0$ (esto es, $T_{\min} = 10\,^{\circ}\mathrm{C}$);
  \item $\partial\Omega_{\mathrm{aislada}}$ es el resto del cerramiento opaco,
  tratado como adiabático (condición de Neumann homogénea).
\end{itemize}

\subsection{Extensión de la definición a condiciones de contorno mixtas}

La Definición~\ref{def:subsuper} está enunciada para el problema de Dirichlet
\eqref{eq:problema_nolineal}, en el que toda la frontera comparte la misma
condición $u = g$. Para trasladarla al problema mixto \eqref{eq:piscina_cc} basta
sustituir, en cada tramo de la frontera, la desigualdad de contorno por la
correspondiente al tipo de condición impuesto en ese tramo:

\begin{itemize}
  \item \textbf{Sub-solución:} $-\Delta \underline{u} \leq f(x,\underline{u})$ en
  $\Omega$; en la puerta $\underline{u} \leq u_{\mathrm{puerta}}$; en la
  cristalera $\partial_n \underline{u} + \alpha\,\underline{u} \leq
  \alpha\,u_{\mathrm{ext}}$; y en el cerramiento adiabático
  $\partial_n \underline{u} \leq 0$.

  \item \textbf{Super-solución:} $-\Delta \overline{u} \geq f(x,\overline{u})$ en
  $\Omega$; en la puerta $\overline{u} \geq u_{\mathrm{puerta}}$; en la
  cristalera $\partial_n \overline{u} + \alpha\,\overline{u} \geq
  \alpha\,u_{\mathrm{ext}}$; y en el cerramiento adiabático
  $\partial_n \overline{u} \geq 0$.
\end{itemize}

Como las tres condiciones de contorno son \textit{lineales} en $u$, la extensión
no altera la demostración del principio de comparación
(Teorema~\ref{teo:comparacion}): ésta sólo recurre al principio del máximo débil
para el operador $L(u) = -\Delta u + M u$ (Teorema~\ref{teo:L_max_debil}), que es
igualmente válido cuando en la frontera se imponen condiciones mixtas de Robin,
Neumann o Dirichlet. Las dos sucesiones monótonas conservan por tanto el
encajonamiento \eqref{eq:monotonía} y la convergencia garantizada por el
Teorema~\ref{teo:existencia}.

\subsection{Parámetros modificables del modelo}

La aplicación interactiva expone cuatro parámetros que alteran la física del
problema. La Tabla~\ref{tab:piscina_parametros} precisa el significado exacto de
cada uno, el símbolo con el que aparece en las ecuaciones
\eqref{eq:piscina_modelo} y \eqref{eq:piscina_cc}, y el rango dentro del cual el
método de sub y super-soluciones está bien definido.

\begin{table}[htbp]
  \centering
  \caption{Parámetros modificables del modelo de la piscina y su rango de
  validez.}\label{tab:piscina_parametros}
  \small
  \begin{tabular}{@{}llllp{4.0cm}@{}}
    \toprule
    \textbf{Parámetro} & \textbf{Símbolo} & \textbf{Rango} & \textbf{Defecto} & \textbf{Significado} \\
    \midrule
    Reacción (evaporación) & $\lambda$ & $\lambda \geq 0$ & $5$ & intensidad de la no
    linealidad logística $\lambda\,u(1-u)$ \\
    Fuente corporal & $\kappa$ & $\kappa \geq 0$ & $2$ & transferencia de calor de cada
    persona ($\kappa_i(x) = \kappa$ en su celda) \\
    Convección (cristalera) & $\alpha$ & $\alpha \geq 0$ & $31.5$ & coeficiente de Robin en
    \eqref{eq:piscina_cc}; acoplamiento adimensional $\gamma = \alpha h$ \\
    Temperatura exterior & $T_{\mathrm{ext}}$ & $T_{\min} \leq T_{\mathrm{ext}}
    \leq T_{p}$ & $18\,^{\circ}\mathrm{C}$ & aire exterior tras la cristalera
    ($u_{\mathrm{ext}} \in [0,1]$) \\
    \bottomrule
  \end{tabular}
\end{table}

Con los valores por defecto y $h = 1/63$ (malla de $64 \times 64$ nodos), el
acoplamiento adimensional de la cristalera vale $\gamma = \alpha h \approx 0.5$.

La restricción sobre $T_{\mathrm{ext}}$ no es arbitraria. La normalización
\eqref{eq:piscina_normalizacion} fija $T_{\min} = 10\,^{\circ}\mathrm{C}$ como la
temperatura más fría del problema, de modo que toda temperatura normalizada
pertenece a $[0,1]$. Si se admite $T_{\mathrm{ext}} < T_{\min}$, entonces
$u_{\mathrm{ext}} < 0$ y el método deja de ser válido por dos motivos: (i) la
constante $\underline{u} \equiv 0$ deja de ser sub-solución, pues en la cristalera
se exigiría $0 \leq \alpha\,u_{\mathrm{ext}}$, que falla; y (ii) la fórmula
discreta de Robin \eqref{eq:piscina_robin_discreta} deja de ser una media convexa
y puede generar valores $u_b < 0$, rompiendo el encajonamiento en $[0,1]$ en el
que se apoya la monotonía de la iteración (Teorema~\ref{teo:comparacion}). Por
esta razón la aplicación restringe $T_{\mathrm{ext}}$ al intervalo
$[T_{\min}, T_{p}]$.

\subsection{Construcción de sub y super-soluciones}

Para aplicar el Teorema~\ref{teo:comparacion} necesitamos un par ordenado de
estimaciones iniciales. En este problema, las constantes son candidatas
naturales y verifican de forma inmediata las desigualdades diferenciales y de
contorno (Definición~\ref{def:subsuper}, extendida a condiciones mixtas):

\begin{itemize}
  \item \textbf{Sub-solución:} $\underline{u} \equiv 0$, el escenario térmico más
  frío posible (piscina vacía, puertas abiertas permanentemente). En el interior,
  \begin{equation*}
    -\Delta \underline{u} = 0 \leq f(x, 0) = \lambda \cdot 0 \cdot 1
    + \sum_{i} \kappa_i(x)\, u_{p,i} \geq 0,
  \end{equation*}
  con igualdad cuando la sala está vacía ($N_{\mathrm{pers}} = 0$). En la
  frontera: en la puerta $\underline{u} = 0 = u_{\mathrm{puerta}}$; en la
  cristalera $0 \leq u_{\mathrm{ext}}$; y en el cerramiento adiabático
  $\partial_n \underline{u} = 0$.

  \item \textbf{Super-solución:} $\overline{u} \equiv 1$, el límite térmico
  superior (aforo completo, cerramiento total). En el interior,
  \begin{equation*}
    -\Delta \overline{u} = 0 \geq f(x, 1) = \lambda \cdot 1 \cdot 0
    + \sum_{i} \kappa_i(x)\,(u_{p,i} - 1) \leq 0,
  \end{equation*}
  pues $u_{p,i} \leq 1$. En la frontera: en la puerta
  $\overline{u} = 1 \geq u_{\mathrm{puerta}} = 0$; en la cristalera
  $1 \geq u_{\mathrm{ext}}$; y en el cerramiento adiabático
  $\partial_n \overline{u} = 0$.
\end{itemize}

Ambas constantes son $C^{2}$ y están ordenadas, $\underline{u} \leq
\overline{u}$, por lo que constituyen un par válido de sub y super-solución.
Obsérvese que la super-solución $\overline{u} \equiv 1$ no iguala el dato de
Dirichlet en la puerta ($u_{\mathrm{puerta}} = 0$): basta con que lo acote por
arriba, tal y como exige la Definición~\ref{def:subsuper}.

\subsection{Parámetro de estabilización $M$}

La derivada parcial de la no linealidad respecto de $u$ es
\begin{equation}\label{eq:piscina_derivada}
  \frac{\partial f}{\partial u}(x,u) = \lambda\bigl(1 - 2u\bigr)
  - \sum_{i=1}^{N_{\mathrm{pers}}} \kappa_i(x),
\end{equation}
cuyo módulo sobre $u \in [0,1]$ está acotado por
$\lambda + \max_{x} \sum_i \kappa_i(x)$. En particular, $f(x,\cdot)$ es
Lipschitz en $u$ sobre $[0,1]$ con constante
$L = \lambda + \max_{x} \sum_i \kappa_i(x)$
(Definición~\ref{def:lipschitz}), y cualquier $M \geq L$ hace que
$F(x,u) = f(x,u) + M u$ sea no decreciente en $u$, que es exactamente la
hipótesis \eqref{eq:F_creciente} del Teorema~\ref{teo:comparacion}. Por la
Proposición~\ref{prop:eleccionM}, el parámetro de estabilización se toma como
\begin{equation}\label{eq:piscina_M}
  M = \lambda + \kappa\, N_{\mathrm{local}},
\end{equation}
donde $N_{\mathrm{local}} = \max_{x} \sum_i \mathbf{1}_{\{\kappa_i(x)>0\}}$ es
el máximo número de personas que comparten una misma celda de la malla. En la
práctica $N_{\mathrm{local}} = 1$ (las personas rara vez ocupan la misma celda),
por lo que con $\lambda = 5$ y $\kappa = 2$ resulta $M = 7$. A diferencia del
modelo de Frank-Kamenetskii, $M$ queda ajustado a la constante de Lipschitz
efectiva, lo que favorece una convergencia rápida.

\subsection{Discretización y estructura de $M$-matriz}

El laplaciano se discretiza mediante la plantilla de cinco puntos sobre una malla
uniforme de $N_x \times N_y$ puntos (en la aplicación interactiva se emplea
$64 \times 64 = 4096$ incógnitas), siguiendo el esquema del
Apéndice~\ref{ape:codigo}. La iteración monótona discreta es
\begin{equation}\label{eq:piscina_iteracion}
  (A_h + M I)\,\mathbf{u}^{(k+1)}
  = \mathbf{f}\bigl(\mathbf{u}^{(k)}\bigr) + M\,\mathbf{u}^{(k)},
\end{equation}
que coincide formalmente con \eqref{eq:iteracion_discreta}. La novedad reside en
el tratamiento de las condiciones de contorno mixtas:

\begin{itemize}
  \item \textbf{Dirichlet (puerta):} el valor $u_{\mathrm{puerta}} = 0$ se fija
  en los nodos de $\partial\Omega_{\mathrm{puerta}}$ cuando la puerta está
  abierta; cuando está cerrada, ese tramo pasa a ser adiabático.

  \item \textbf{Robin (cristalera):} discretizando la derivada normal con una
  diferencia finita descentrada de paso $h$, la condición
  \eqref{eq:piscina_cc} sobre un nodo frontera $u_b$ con vecino interior
  $u_{\mathrm{int}}$ se reduce a
  \begin{equation}\label{eq:piscina_robin_discreta}
    u_b = \frac{u_{\mathrm{int}} + \gamma\, u_{\mathrm{ext}}}{1 + \gamma},
    \qquad \gamma = \alpha\, h,
  \end{equation}
  donde $\gamma$ es el acoplamiento convectivo adimensional. Nótese que la
  fórmula es una media convexa de $u_{\mathrm{int}}$ y $u_{\mathrm{ext}}$, lo que
  respeta el principio del máximo.

  \item \textbf{Neumann (cerramiento):} se impone $u_b = u_{\mathrm{int}}$ en
  los nodos adiabáticos (espejo del vecino interior).
\end{itemize}

Con independencia de la condición aplicada, la matriz del sistema $A_h + M I$
sigue siendo una \textbf{$M$-matriz} no singular. Las filas interiores son
diagonalmente dominantes estrictas (diagonal $4/h^{2} + M$ positiva y
extradiagonales $-1/h^{2}$ no positivas); las filas de Robin tienen diagonal
$1 + \gamma$ y única extradiagonal $-1$, también dominantes estrictas para
$\gamma > 0$; y las filas de Neumann ($u_b = u_{\mathrm{int}}$) son sólo
dominantes en sentido débil. No obstante, al ser el dominio conexo la matriz es
\textit{irreduciblemente} diagonalmente dominante, condición suficiente para que
sea una $M$-matriz no singular. En consecuencia $(A_h + M I)^{-1} \geq 0$ y el
principio del máximo discreto se preserva (Observación~\ref{obs:max_discreto}),
lo que garantiza que las dos sucesiones generadas desde $\underline{u}$ y
$\overline{u}$ sean monótonas, acotadas y converjan a la distribución real de
temperaturas del edificio.

\subsection{Resultados numéricos y visualización interactiva}

La Figura~\ref{fig:piscina_sub_super} muestra la convergencia del método. La
sub-solución (partiendo de $u \equiv 0$) \textbf{crece} monótonamente mientras
la super-solución (partiendo de $u \equiv 1$) \textbf{decrece}, encajonando la
solución de equilibrio $u^{*}$ en la cadena de orden
$\underline{u} \leq u_1 \leq \cdots \leq u^{*} \leq \cdots \leq \overline{u}_1
\leq \overline{u}$ predicha por el Teorema~\ref{teo:comparacion}.

\begin{figure}[htbp]
  \centering
  \includegraphics[width=0.85\textwidth]{piscina_sub_super.png}
  \caption{Convergencia del método de sub y super-soluciones para el modelo de
  climatización de la piscina. \textbf{Arriba:} sub-solución inicial ($u \equiv
  0$, escenario más frío) y super-solución inicial ($u \equiv 1$, aforo
  completo). \textbf{Abajo:} solución de equilibrio con las fuentes corporales
  (personas) y el frente frío de la puerta abierta.}
  \label{fig:piscina_sub_super}
\end{figure}

El perfil de equilibrio resultante refleja los dos mecanismos del modelo: un
frente frío que penetra desde la puerta abierta (condición de Dirichlet) y una
capa templada junto a la cristalera (condición de Robin), mientras las personas
generan focos locales de calor que se difunden por la sala. El decaimiento del
\textit{gap} $\|\overline{u}^{(k)} - \underline{u}^{(k)}\|_{\infty}$ en escala
logarítmica (Figura~\ref{fig:piscina_convergencia}) confirma la convergencia
lineal asintótica característica del método.

De acuerdo con el Teorema~\ref{teo:existencia}, las dos sucesiones convergen en
$C^{2}$ a sendas soluciones $u_{\min}$ y $u_{\max}$ del problema mixto, con el
encajonamiento $\underline{u} \leq u_{\min} \leq u^{*} \leq u_{\max} \leq
\overline{u}$ (\eqref{eq:encajonamiento}). Puesto que $f$ es cóncava en $u$
($\partial_{u}^{2} f = -2\lambda < 0$), la unicidad en el intervalo $[0,1]$ está
sugerida (Corolario~\ref{coro:unicidad}); numéricamente ambas sucesiones
convergen al mismo perfil de equilibrio, con el encajonamiento bilateral actuando
como certificado de precisión.

\begin{figure}[htbp]
  \centering
  \includegraphics[width=0.85\textwidth]{piscina_convergencia.png}
  \caption{Decaimiento del error en escala logarítmica para el modelo de la
  piscina. El \textit{gap} $\|\overline{u}^{(k)} - \underline{u}^{(k)}\|_{\infty}$
  entre ambas sucesiones decrece de forma aproximadamente lineal en escala
  semilogarítmica, evidenciando la convergencia monótona hacia la solución de
  equilibrio.}
  \label{fig:piscina_convergencia}
\end{figure}

A diferencia de los ejemplos 1D, este modelo cuenta con una \textbf{aplicación
3D interactiva} (desarrollada con TypeScript y Three.js) que permite añadir o
retirar personas en tiempo real y observar cómo ambas sucesiones se reajustan
hacia el nuevo equilibrio, así como alternar la visualización entre la
sub-solución, la super-solución, el \textit{gap} $\overline{u} - \underline{u}$
y la solución de equilibrio. La Figura~\ref{fig:piscina_app} muestra una captura
de dicha aplicación, que complementa la validación numérica con una
visualización orientada a la defensa del trabajo.

\begin{figure}[htbp]
  \centering
  \includegraphics[width=0.9\textwidth]{piscina_app.png}
  \caption{Captura de la aplicación 3D interactiva del modelo de la piscina. El
  mapa de calor sobre el suelo representa la temperatura normalizada; los
  controles permiten modificar el aforo, abrir o cerrar la puerta y alternar
  entre sub-solución, super-solución y solución de equilibrio.}
  \label{fig:piscina_app}
\end{figure}

% ============================================================================
%  OPCIONAL: actualización de la Tabla~\ref{tab:contraste} ("Contraste de los
%  Modelos") para incluir este cuarto ejemplo. Descomentar y añadir una cuarta
%  columna a dicha tabla, junto con las filas correspondientes.
% ============================================================================
%
% En el preámbulo de la tabla, cambiar
%   \begin{tabular}{@{}lccc@{}}
% por
%   \begin{tabular}{@{}lcccc@{}}
% y en la cabecera añadir "& \textbf{Piscina 2D} \\".
%
% Filas a añadir (alineadas con las ya existentes):
%   $f(u)$                              & $\lambda e^{u}$   & $\lambda u(1-u)$   & $\lambda e^{u}$ & $\lambda u(1-u) + \sum_i \kappa_i(x)(u_{p,i}-u)$ \\
%   No linealidad                       & Exponencial       & Logística          & Exponencial     & Logística + fuentes locales \\
%   Convexidad                          & Convexa           & Cóncava            & Convexa         & Cóncava en $u$ (depende de $x$) \\
%   Dimensión                           & 1D                & 1D                 & 2D              & 2D \\
%   Incógnitas                          & 100               & 100                & 2\,500          & 4\,096 \\
%   $L$ (Lipschitz)                     & $\lambda e^{u_{\max}}$ & $\lambda$     & $\lambda e^{u_{\max}}$ & $\lambda + \kappa N_{\mathrm{local}}$ \\
%   $M$ empleado                        & $\lambda e^{2.5} \approx 24.36$ & $6$ & $\lambda e^{2.5} \approx 12.18$ & $7$ \\
%   $\lambda$                           & $2$               & $5$                & $1$             & $5$ \\
%   Cond. Contorno                      & $u=0$ extremos    & $0.9,\,0.1$ extremos & $u=0$ en $\partial\Omega$ & Robin + Dirichlet + Neumann \\
%   Iteraciones (sub)                   & $48$              & $12$               & $14$            & (medido en la app) \\
%   Iteraciones (super)                 & $53$              & $12$               & $17$            & (medido en la app) \\
%   Unicidad                           & No trivial\textsuperscript{1} & Sugerida & No trivial & Sugerida \\
